% 1. Une discussion sur à quel point la station actuel est mauvaise

\subsection{Capteurs actuels}
Les capteutrs actuel se résument à une qualité d'amateurs.
Par exemple, le DHT11 à des températures de fonction entre 0 et 50 degrées celcius, ce qui ne marchera pas l'hivers.
De plus, il lit l'humidité que de 20\% à 90\%, avec une stabilité médiocre entre capteurs. Il ne peut donc pas être utilisé.

Le capteur DSP310 est un bon capteur. Il fonctionne de -40C à 85C, à une précision de $\pm$ 0.5C et $\pm$ 1 hPa.
Il est plutôt fidèle et fiable, et est utilisé dans des situations commerciales. Sa consommation est optimal, à environ 2$\mu$A.
Il est donc correct de continuer à l'utilisé.

Pour ce qui est de l'ensolleillement, le LS06-S fonctionne de -25C à 75C, donc quasiment correct pour nos hivers. Sa consommation
dépend de la luminosité extérieur, autour de quelques mA. Cependant, son gros problème est le manque d'information dans sa datasheet
pour être en mesure de convertir de façon commercial, ses valeurs en unités connues. Ayant donc une fidélité médiocre.
Il fonctionne mieux comme un capteur de lumière que quelquechose qui donne une valeur d'ensolleillement précis.
De plus, il sature facilement.

Le pire reste cependant la station météo sparkfun elle même. C'est un kit éducationel qui n'est pas fait pour une usage commerciale.
En effet, sa composition en plastique est fébrile et ne supporterais pas d'être à l'extérieur. La présence de températures en dessous
de 0C vont causé des bris avec le gel. La précision de l'anémomètre est au meilleur, médiocre. La girouette pour la direction du vent
à des résistances tellement pas précise qu'il requiert une calibration des directions différentes pour chaque unitées

\subsection{Protocoles actuels}

\subsection{Microcontrolleur actuel}